% !TeX root = ../main.tex
\section{Current results and their boundaries}\label{sec:zk-status}

The following are the main entry points, not pieces of an already composable construction. Their bounds refer to different views and, sometimes, different padding layouts or setup models.

\paragraph{Immediate obstruction: frozen count projections.} \noteref{zk-count-frontier} proves that all $129788$ balanced-source metadata bits preserve every aligned $2^{14}$-leaf count product; compression-input and unused-memory masks do too. A private coarse frontier is exposed by the sixth GKR packet except with probability at most $10/2^{192}$, despite equal products in all twenty-eight count columns. Localization can make that frontier public, subject to an unproved capacity contract, but does not fix the whole count head: a second valid private-pointer pair inside one $256$-row interval still leaks through packet eleven's count child zero. These obstructions reject specific fixed completions, not the general padding-only idea.

\paragraph{A targeted in-frame count repair.} \noteref{zk-count-adapter} adds $3072$ MUL adapters to existing compression cycles. Their $1536$ independent label trades reach that later count child while preserving complete read chains, all column products and compression witnesses. No compression rows or frame allocations are added. The child has a uniform conditional span outside an event below $2^{-160}$, retaining all $32$ newly exposed prefix bits. It joins the established restricted boundary below $2^{-151}$, with a $1116$-bit law for the entire $4096$-point prefix. Earlier GKR messages and the new MUL disclosures are excluded. The adapters preserve the coarse frontier, so its repair, global reservation and the all-layer joint proof remain necessary.

\paragraph{Two query regions and an entire public prefix.} Theorem~\ref{thm:zk-prefix-split-interface} in \noteref{zk-flock-lowbank} jointly hides one sixteen-point coset in $\subsp_{18}\setminus\subsp_{13}$ and one thirty-two-point coset in $\gen^{18}+\subsp_{18}$ while retaining all $4096$ values on $\subsp_{12}$ with their actual correlated law. The view includes the twelve final GKR children, all BLAKE2s terminal claims, established Flock endpoints, memory envelope and shared root, below $2^{-151}$. The same $3840$ cycles suffice: independent count columns separate the regions, an old count bank hides the terminal claim before the new raw-query masks, and $188$ visible PC bits are retained explicitly. The prefix encodes $1084$ independent source bits. Revealing the entire larger prefix $\subsp_{13}$ would instead leak through six further values (Proposition~\ref{prop:zk-full-prefix-envelope}); this is an overdisclosure obstruction, not an attack on the honest-query target. Common valid completion and the public-prefix contract remain assumed; arbitrary queries and the complete opening remain unproved.

\paragraph{Why the bank and matching matter.} Without the bank, \noteref{zk-flock-query-fiber} proves that six low queries distinguish two valid private-pointer witnesses, with common-simulator error floors above $2^{-93},2^{-105},2^{-114},2^{-123}$ across the four rates. The original long-chain matching repaired that test but retained a twenty-eight-query dependency and a structural thirty-two-point rank deficit. The balanced matching removes those specific deficits using the same rows. The complete exceptional-query family remains unbounded.

\paragraph{A genuine dispersed dependency.} Proposition~\ref{prop:zk-scattered-support} certifies a minimal raw distinguishing support of $42$ points for the balanced source, one per sixteen-point coset. The whole noise code on $\subsp_{14}\setminus\subsp_{13}$ has binary dimension $2688$; this particular projection has rank $2687$, with every proper subset full rank. The probability of at least $42$ queries anywhere in that band is below $2^{-266}$ at rate one. Thus the displayed dependency is affordable, but a proof based only on concentration within small cosets is insufficient. Smaller and cross-band dependencies, particularly after earlier disclosures, remain unbounded.

\paragraph{Memory: an obstruction and a restricted-program repair.} \noteref{zk-memory-frames} exhibits a pair of complete valid executions for which unrestricted execution-preserving padding cannot reach high statistical privacy: low-point PCS queries expose fixed private memory prefixes. A public small-frame allocation contract avoids this obstruction. Its pin-preserving construction realizes unused-memory masks and hides a complete algebraic memory envelope, including later folds. \noteref{zk-memory-root} composes that envelope with the shared bus root below $2^{-151}$, without further reserved words. The root-replacement hybrid also retains the actual mixed algebraic PCS and other non-GKR reductions, but does not simulate their actual marginal. Intermediate fingerprint GKR messages, joint simulation of the other columns and authentication remain open. This is the current P1/P5 feasibility branch, not full ZK.

\paragraph{Ordinary range checks can stay unchanged.} Theorem~\ref{thm:zk-public-probe-envelope} keeps their possible low-address probes public and moves the unused segment-zero masks beyond that prefix. Corollary~\ref{cor:zk-probe-relocated-interface} carries the established restricted BLAKE2s/count-query interface through the frame relocation, jointly with memory and the shared root below $2^{-151}$. At the candidate height this costs $256$ free slots, leaving $26703$ before fragmentation, with no new range-check rows or verifier equations. Native proofs validate a public bootstrap and the unchanged gadget; the full metadata-cycle audit checks the relocated masks and boundary. General compilation, total allocation and full transcript remain unproved.

\paragraph{The actual guest's code height is compatible with the restricted proof.} Corollary~\ref{cor:zk-large-code-interface} extends that interface to bytecode log $19$, retaining the bound below $2^{-151}$. The changed PCS placement requires different low/high count columns and two replacement GKR minors; both old minors become identically zero. The complete $4096$-point prefix now encodes $1085$ fair bits. All $99$ padding codes fit in the native guest's unused tail after a fixed public translation. The native code-space check and algebraic theorem do not install the new bytecode or prove the general compiler/allocator transformation. Earlier messages, other lanes and queries, authentication and full ZK remain open.

\paragraph{An executable allocator and the wrapped aggregation guest.} Lemma~\ref{lem:zk-public-entry-wrapper} and Proposition~\ref{prop:zk-reserved-allocation} give a public entry construction and a relocation argument under an explicit address discipline and total-allocation bound. The \texttt{zk-research} feature routes static/runtime allocation hints and filler frames through the interval allocator. Besides the original private-branch fixture, the actual wrapped aggregation guest passes its code/read-cap checks and leaf execution audits for both signature schemes. A $453$-XMSS witness fits the reserved memory and private compression budget; a smaller wrapped leaf has a valid native proof (\S\ref{par:zk-guest-allocation-audit}). The tests preserve unused masks, corresponding frame/instruction identities and compression data. Default allocation and all verifier equations are unchanged. The address discipline and resource bound still need certification for every supported private execution, including recursive children, followed by the integrated ZK cycle sampler. This is a P1 mechanism and conditional legality theorem, not full ZK.

\paragraph{The public probe prefix still has private counts.} Proposition~\ref{prop:zk-probe-count-floor} identifies a separate uncovered lane: at the candidate layout, query zero in lane $40$ reveals the final read count of memory cell zero. Keeping only its value public leaves simulator-error floors above $2^{-17}$ through $2^{-22}$ across the four rates when padding adds the same fixed count. Lemma~\ref{lem:zk-probe-count-completion} repairs a protected count prefix by fixed-cost real/dummy read cycles with public quotas. Four native proofs certify the leaking pair and a $128$-cell completion at one common statement; the complete cycle library is checked separately. The repair transfers private information into helper and instruction counts, and its full-prefix cost does not automatically fit the candidate. This rejects the value-only completion, not the small-frame approach or the restricted theorems, which excluded this count lane.

\paragraph{The prefix repair can preserve public column products.} Proposition~\ref{prop:zk-probe-column-completion} jointly fixes the protected final counters and all twenty-eight count-column products using only MUL, DEREF and JUMP additions, assuming XOR, SET and compression products are already public. Corollary~\ref{cor:zk-noncompression-completion} also handles initially private XOR/SET products without adding or relabeling compression rows. Exact cycle audits close the exhibited residual root leak, including all $27$ bounded three-cell access patterns with privately varying totals. Corollary~\ref{cor:zk-probe-count-memory-root} joins these public count coordinates to the memory/root marginal below $2^{-151}$. The explicit row ledger is conditional on public exponent intervals and disjoint capacity; it does not establish candidate-wide feasibility, normalize the remaining eighteen coarse frontiers, or hide intervening GKR and other count queries.

\paragraph{Small frames alone are insufficient at a fixed probe budget.} Corollary~\ref{cor:zk-probe-budget-query} gives an actual honest-query obstruction for every randomized execution-preserving completion under the current outside-prefix frame policy: alternative private read loads above half the available prefix-access budget force disjoint count-zero supports. Native five-read proofs and $262145$-read executions certify the mechanism and code capacity. The program class needs a public probe-load restriction or a different prefix-access policy. Separately, Proposition~\ref{prop:zk-count-prefix-budget} rules out the stronger full-count-prefix envelope with error below $1/9$ for a nine-range-check family at the candidate budget. That second bound is not an attack on the actual sparse-query verifier. Neither result rules out a capacity-restricted aggregation construction.

\paragraph{Canonical leaves need only eight private-count quotas.} \noteref{zk-leaf-probes} identifies a more economical public contract for raw aggregation leaves with externally public signer geometry. Coverage bijections fix their entire probe histogram; other metadata probes are public. A bytecode-only local DEREF filler removes the private filler contribution at index zero without new rows. Only the eight WOTS digit addresses remain private. Tight target-sum quotas equalize them with eight shared helper frames and exactly $230$ DEREF/JUMP cycles per encoding word, making all $65536$ prefix counters public. Complete cycle certificates and native leaf profiles support this result; a native wrapped proof verifies the optional filler rewrite. Helper/root normalization, common capacity, recursive children and the full joint view remain open. The geometry must genuinely be available to the simulator, not inferred from its digest.

An optional local polynomial check can instead remove all eight digit probes (Proposition~\ref{prop:zk-local-digit-membership}). Its membership relation is exact over $\E$, so all canonical-leaf prefix counters become public without count top-ups. Native mixed proofs and the $453$-XMSS allocation check pass; invalid-input regressions exercise the compiled check. This is a public guest rewrite with additional arithmetic rows, not free padding or a change to the three verifier algorithms. Full arithmetic/count normalization, recursive probes and joint statistical ZK remain open.

\paragraph{All range probes can be localized.} \noteref{zk-local-ranges} replaces all $70$ range-check sites by local Boolean/exponent certificates, with logarithmic cost in each bound and no new verifier equations. The exact range lemma includes private runtime bounds under their existing cap obligation. With a public wrapper and valid object ownership, prefix counts now depend only on the fixed bootstrap and two input reads, not on private child geometry. Mixed-leaf and one- and two-child outer proofs verify; the larger leaf still fits the necessary allocation and opcode budgets. Two masked-read tails needed explicit allocation in the optional program. Common padding, a universal recursive resource/ownership theorem, the remaining count frontiers and full joint statistical ZK remain open.

\paragraph{Final GKR layer: a joint four-round tail.} \noteref{zk-four-round-tail} hides the last four sumcheck rounds of the final GKR layer, their incoming claim, final children and five framework evaluations, below $2^{-152}$. It uses forty-eight sparse row-swap banks, respects native push/pull/count relations and retains a private self-loop invariant. The supported layer has twenty rounds: its first sixteen rounds are excluded, as is composition with earlier layers. The theorem assumes compatible count normalization and a disjoint valid completion. It does not include the subsequent table proof or PCS.

\paragraph{Coarse GKR layers: a separate setup branch.} \noteref{zk-mul-control} jointly simulates the first two GKR layers and the seventeen-field terminal boundary below $2^{-155}$, outside an honest public-setup exception below $2^{-222}$. It requires public random word/frame-scale libraries and native field-valued frame/offset support that is not implemented. These are additional, unadopted assumptions. The intervening layers are omitted; this theorem cannot simply be joined to the four-round tail theorem.

\paragraph{Counts: normalization and a partial internal simulator.} \noteref{zk-count-columns} makes all twenty-eight count-column products public using valid cycles and label routing, subject to capacity and compatibility conditions. This does not hide the internal count tree. \noteref{zk-count-uniform-tail} proves an eight-round count-only tail theorem below $2^{-157}$. \noteref{zk-count-completed-head} supplies exact conditional head-rank certificates for that construction, but not a uniform head failure bound or a joint simulator with the fingerprint channels.

\paragraph{Small-frame count completion.} Proposition~\ref{prop:zk-reused-count-frames} removes the normalization stages' per-traversal frame allocations. Forty reused frames normalize noncompression bytecode; the resulting memory uncertainty has an exact affine formula and is handled by the subsequent stages. The whole conditional all-column completion needs $62$ frames of at most $128$ cells, with compression rows unchanged. Sixteen complete cycle certificates check independent variations of code multiplicity and memory reuse. Public opcode totals, updated exponent intervals and sufficient row/code capacity are still premises; coarse-frontier hiding and a common candidate sampler are not established.

\paragraph{A constructive route for all remaining coarse products.} \noteref{zk-count-coarse-balance} gives a conditional simultaneous repair of all $704$ noncompression coarse nodes using balanced row order and independent count-label transfers. Shared XOR/MUL/SET/DEREF cycles halve the coarse router's JUMP cost to $221184$, with $96$ frames and $224$ instructions at label cap $165888$. Sixteen-instruction fillers close the residual row ledger exactly, including every return; Corollary~\ref{cor:zk-common-coarse-budget} states the sufficient and necessary inequalities for this completion stage. Four integrated cycle certificates run whole-column normalization and coarse balancing together, preserving all compression rows and agreeing in final products, native height-seven nodes and memory/code images despite private code/memory variation. Fixed masking backdrops are checked separately. The guest-wide incoming public totals, postnormalization label cap and common reservation still need proof. Public coarse products alone do not simulate the full joint GKR transcript.

\paragraph{A smaller, correlated incoming ledger.} \noteref{zk-count-budget} cancels JUMP's bytecode uncertainty and common memory shifts from its later router widths. Unequal centers and rectangular routers give a large-interval input contract that fits all five tables, every filler return and the label cap; the previous square-router ledger overfills MUL for that same contract. Complete cycle certificates and native coarse-frontier compositions check both variants. Local-role priority additionally bounds the two original JUMP differences from the actual public bytecode, independently of foreign indirect reads. Under the fresh-frame contract, their widths fit the numerical contract through $46136$ real JUMPs; the initial fillers preserve them. Native label reconstruction checks every real final counter and the conserved total. The remaining count ranges, uniform real-row budgets, initial public totals, full allocation and joint statistical privacy remain unproved. This is not yet a common guest-wide sampler.

\paragraph{The BLAKE2s coarse frontier: a conditional repair.} Theorem~\ref{thm:zk-bc-coarse-normalization} fixes all sixteen bytecode-count products for the $65536$-row private budget, reusing $95112$ mandatory rows. Theorem~\ref{thm:zk-memory-coarse-normalization} and Lemma~\ref{lem:zk-group-label-router} also fix the $144$ memory-count products under a public per-cell load bound. The stronger router supports bound $256$ with $512$ reused compressions, $149319$ XOR rows and $16591$ extra returns. Relocating one already-erased low-bank block preserves the restricted theorem below $2^{-151}$, with $1117$ visible prefix bits. The actual aggregation bytecode passes the static read-cap audit under the fresh-frame contract; large frames can use consecutive free slots. Code reservation, the general compiler/allocator transformation and total allocation still prevent claiming a common sampler. The remaining eighteen count columns, transferred XOR/JUMP uncertainty and full count head remain open.

\paragraph{Table constraints: a conditional composition route.} \noteref{zk-jump-rank} hides the isolated JUMP constraint sumcheck; \noteref{zk-jump-metadata} adds its terminal columns. \noteref{zk-table-batch} reduces simulation of the actual six-table sumcheck to a larger preceding view that still needs simulation. \noteref{zk-jump-bus-padding} hides the bus-batched JUMP component and its boundary below $2^{-146}$, but does not establish its joint privacy with preceding GKR messages.

\paragraph{Flock, twelve GKR children and twelve lane values.} \noteref{zk-flock-sixteen} jointly hides all thirty-seven BLAKE2s terminal claims, the established Flock endpoints, all twelve transmitted final GKR children and any twelve lane-$59$ values in one sixteen-point encoder coset, below $2^{-155}$. The unchanged library now supplies eight kernel and four quotient coordinates. A weighted cubic dual proves that every twelve-point projection discards its four surviving invariants. Shared child/terminal/query bits are covered by dense fixed-disclosure budgets of $128$ and $320$ bits. The theorem concerns one coset inside $\gen^{18}+\subsp_{18}$, not independent simulations for many cosets. Preceding GKR rounds, table coefficients, middle Flock rounds and common allocation remain open.

\paragraph{A cross-subsystem composition.} Theorem~\ref{thm:zk-flock-twelve-interface} also joins that enlarged boundary to the shared bus root and complete three-limb memory envelope below $2^{-151}$. The proof first replaces metadata, children and the selected query values, then replaces memory/root while retaining those uniforms, and finally simulates the compression-value/Flock interface. All actual queries in one selected sixteen-point coset may be included at the same stated exponent by charging the event of thirteen distinct queries there, below $2^{-174}$ at rate one. This is not a simulation of all queries. The theorem assumes a valid common completion within the public reservation; the count root, preceding GKR messages, other PCS values and authentication are excluded.

\paragraph{PCS: restricted algebraic openings.} \noteref{zk-pcs-padding} hides a complete algebraic WHIR opening for one independent random-point claim, assuming freely random padding coefficients. \noteref{zk-stacked-padding} extends this to one common Frobenius orbit. Neither covers the full actual stacked claim family jointly with the VM reductions, nor realizes every required coefficient with valid cycles. Commitments and Fiat-Shamir are excluded.

\paragraph{Authentication feasibility.} \noteref{zk-memory-leaf-entropy} proves that ordinary unopened initial leaves retain at least $384$ bits of conditional entropy after the memory envelope, outside a common event below $2^{-161}$. The averaged probability of any coinciding initial input images is below $2^{-909}$. The entropy statement does not yet survive conditioning on the full GKR transcript and is not a hash-simulation theorem. The two low encoder points remain a separate opening obligation.


\subsection{How the local results may be used}\label{sec:zk-composition}

\paragraph{One joint view.} The target contains the commitments, announced sizes, bus and count products, every GKR message, table sumcheck, Flock zerocheck and lincheck, ring switching, the complete stacked PCS opening, its final plaintext table, authentication paths and externally visible deferred claims. The protocol order is specified in \path{doc/leanvm/body/08-end-to-end-protocol.tex}; the proof inventory must also be checked against the implementations. No later disclosure of a padding value may be silently omitted from its privacy analysis.

Individual hiding claims do not compose automatically. If a fair bit $U$ masks a secret bit $b$, both $U$ and $U\oplus b$ are individually uniform, but their pair reveals $b$. Likewise, two sumchecks can each look random while revealing the witness through their shared masks. What is needed is a joint simulator, or conditional hybrids that retain every already exposed coordinate and justify the remaining randomness. Only then may their statistical errors be added.

\paragraph{The useful proof pattern.} Identify the native affine relations between the observations, then prove that legal padding choices make the remaining coordinates uniform on the appropriate affine set, except on explicitly bounded exceptional coins. Factor that distribution conditionally to expose unused randomness, and analyze the nonlinear messages it hides. A witness-dependent invariant need not itself become uniform if the resulting observable distribution is independent of it. But it cannot be passed to the final simulator as an unexplained input.

\paragraph{Evidence standards.} An exact nonzero determinant can certify a nonzero polynomial; a degree bound and the actual challenge distribution can then bound its failure probability. A successful numerical rank at one point, especially a large binary rank, is not such a bound. An audit replaying selected messages through the reference verifier checks their algebra, not the omitted messages or a full VM execution. Our randomness arguments use genuine random choices; a pseudorandomness assumption alone does not establish statistical privacy.


\paragraph{Bottom line.} The existing bounds show that native-field statistical errors well below $2^{-128}$ are possible for selected projections. They do not show that the remaining messages have sufficient legal randomness, or that all reservations fit together. A fifth tail round would be progress on one obligation, not a route around the missing full-layer, PCS and authentication theorems.
